\section{Estado actual}

La versión \versionTesis{} convierte la primera entrega en un borrador de tesis consistente. Conserva su problema central --la necesidad de una infraestructura reproducible para estudiar ecosistemas CEX/DEX-- y delimita mejor el artefacto académico: un entorno de simulación y evaluación, no una promesa de rentabilidad ni un sistema de operación financiera.

En este punto se han definido la pregunta, los objetivos, los agentes mínimos, las capas de la arquitectura, los escenarios iniciales y los criterios de evaluación. El repositorio asociado establece el espacio para documentación, notebooks, datos, código y el manuscrito fuente. Aún no se reportan experimentos empíricos ni comparaciones de desempeño.

\section{Hoja de ruta}

La evolución propuesta se organiza en cuatro incrementos:

\begin{enumerate}
  \item \textbf{Fundación de datos.} Seleccionar fuentes iniciales, documentar licencia y procedencia, implementar validación de esquema y publicar un notebook de exploración.
  \item \textbf{Simulación mínima.} Implementar el AMM o libro de órdenes simplificado, los agentes básicos y un escenario de shock reproducible.
  \item \textbf{Evaluación.} Añadir costos de ejecución, baselines, sensibilidad de parámetros y reporte estructurado de métricas.
  \item \textbf{Extensión.} Incorporar agentes adaptativos, componentes multi-resolución, grafos de exposición y escenarios adicionales solo cuando los incrementos anteriores sean estables.
\end{enumerate}

Cada incremento deberá generar una versión de manuscrito, un conjunto de artefactos ejecutables y una nota de cambios. Las afirmaciones del capítulo de resultados se redactarán después de ejecutar ese ciclo, con evidencia enlazada y limitaciones explícitas.

\section{Limitaciones y consideraciones éticas}

Una simulación no reproduce por completo un mercado real. Los resultados dependerán de las fuentes disponibles, las reglas de los agentes, el modelo de costos, el intervalo temporal y las simplificaciones de protocolo. En especial, la representación de MEV, liquidaciones y riesgo sistémico debe interpretarse como un modelo experimental y no como una reconstrucción exhaustiva de una cadena o venue particular.

El proyecto no administrará fondos de terceros, no ejecutará operaciones reales con capital significativo y no ofrecerá asesoría de inversión, jurídica o tributaria. Los datos y artefactos públicos se publicarán respetando licencias, privacidad y restricciones de las fuentes. La transparencia de supuestos será tan importante como cualquier métrica de desempeño.

\section{Cierre}

El valor esperado de la tesis no reside en predecir el próximo movimiento de un activo, sino en ofrecer un instrumento para examinar cómo interactúan liquidez, reglas de ejecución, incentivos y riesgo en mercados digitales híbridos. Al mantener el manuscrito, el código, los notebooks y la evidencia dentro de un mismo flujo versionado, la investigación puede ser discutida, auditada y ampliada como un objeto reproducible.
